\chapter{Integral Dependence}
\label{ch:v19}

Integral dependence generalizes algebraicity from field extensions to ring extensions. An element is integral when it satisfies a monic polynomial, and integrality is stable under finite algebraic constructions.

\section*{Learning goals}
After this chapter, the reader should be able to:
\begin{itemize}
\item test whether an element is integral over a subring by finding a monic polynomial relation;
\item prove closure of integral elements under addition and multiplication using finite-module criteria;
\item apply the determinant trick or finite-generation criterion for integral dependence;
\item analyze integral extensions under quotients and localization;
\item show how integrality constrains prime ideals through lying-over-type statements developed in the chapter;
\item distinguish integral dependence from algebraic dependence or arbitrary polynomial relations;
\end{itemize}
\section*{Conceptual roadmap}
\[
\boxed{\text{presentation or universal property}\;\longrightarrow\;\text{exactness/localization}\;\longrightarrow\;\text{structural consequence}.}
\]

\section{Integral elements}
An element $b$ of an $A$-algebra $B$ is integral over $A$ if it satisfies a monic polynomial with coefficients in $A$.

\section{Integral extensions}
An extension $A\subset B$ is integral when every element of $B$ is integral over $A$.

% BEGIN VOL05-EXPANSION V19-example-01
\begin{example}[Integral element over a subring]\label{ex:v19-expand-01}
Let $A=k[t^2,t^3]\subset k[t]$. The element $t$ is integral over $A$ because it satisfies the monic equation $X^2-t^2=0$ with coefficient $t^2\in A$.
\end{example}
% END VOL05-EXPANSION V19-example-01
\section{Finite module criterion}
An element is integral iff the subalgebra it generates is finite as an $A$-module.

\section{Finite algebras}
Every finite $A$-algebra is integral over $A$.

\section{Transitivity}
Integrality is transitive through towers of ring extensions.

% BEGIN VOL05-EXPANSION V19-example-02
\begin{example}[Gaussian integers]\label{ex:v19-expand-02}
The element $i$ is integral over $\mathbb Z$ because it satisfies $X^2+1=0$. Consequently every element of $\mathbb Z[i]$ is integral over $\mathbb Z$.
\end{example}
% END VOL05-EXPANSION V19-example-02
\section{Integral elements form a ring}
The set of elements of $B$ integral over $A$ is a subring.

\section{Determinant trick}
Finite-generation plus stability under multiplication produces monic equations by a determinant argument.

% BEGIN VOL05-EXPANSION V19-example-03
\begin{example}[Finite-module criterion]\label{ex:v19-expand-03}
If $b$ is integral over $A$, then $A[b]$ is generated as an $A$-module by $1,b,\ldots,b^{n-1}$ for a monic equation of degree $n$. Integrality converts an algebra extension into finite module data.
\end{example}
% END VOL05-EXPANSION V19-example-03
\section{Prime behavior preview}
Integral extensions satisfy lying-over and contraction properties for prime ideals.

\section{Examples}
Algebraic integers, roots of monic equations, and normalization rings illustrate integral dependence.

\section{Core structural results}

\begin{theorem}[Finite module criterion]\label{thm:v19-01}
For $b\in B$, $b$ is integral over $A$ iff $A[b]$ is finitely generated as an $A$-module.
\begin{proof}
A monic equation expresses every high power of $b$ in terms of lower powers. Conversely, if finitely many powers generate an $A$-module stable under multiplication by $b$, the determinant trick gives a monic polynomial annihilating $b$.
\end{proof}
\end{theorem}

\begin{theorem}[Finite algebra implies integral]\label{thm:v19-02}
If $B$ is a finite $A$-module and an $A$-algebra, then every $b\in B$ is integral over $A$.
\begin{proof}
Multiplication by $b$ is an $A$-linear endomorphism of the finite module $B$. Apply the determinant trick or Cayley--Hamilton to obtain a monic polynomial with coefficients in $A$ annihilating $b$.
\end{proof}
\end{theorem}

\begin{theorem}[Transitivity of integrality]\label{thm:v19-03}
If $B$ is integral over $A$ and $C$ integral over $B$, then $C$ is integral over $A$.
\begin{proof}
For $c\in C$, its monic equation uses finitely many coefficients from $B$. Those coefficients generate a finite $A$-algebra because they are integral; adjoining $c$ is finite over that finite algebra, hence finite over $A$, so $c$ is integral over $A$.
\end{proof}
\end{theorem}

\section{Worked examples}

\begin{example}[Integral elements diagnostic]\label{ex:v19-worked-01}
Give a rigorous worked analysis of Integral elements. State the ring, module, finiteness, localization, or homological hypotheses and derive the central conclusion. An element $b$ of an $A$-algebra $B$ is integral over $A$ if it satisfies a monic polynomial with coefficients in $A$. The decisive step is to use the universal property, exactness criterion, localization test, or finite-generation mechanism appropriate to the algebraic structure.
\end{example}

\begin{example}[Integral extensions diagnostic]\label{ex:v19-worked-02}
Give a rigorous worked analysis of Integral extensions. State the ring, module, finiteness, localization, or homological hypotheses and derive the central conclusion. An extension $A\subset B$ is integral when every element of $B$ is integral over $A$. The decisive step is to use the universal property, exactness criterion, localization test, or finite-generation mechanism appropriate to the algebraic structure.
\end{example}

\begin{example}[Finite module criterion diagnostic]\label{ex:v19-worked-03}
Give a rigorous worked analysis of Finite module criterion. State the ring, module, finiteness, localization, or homological hypotheses and derive the central conclusion. An element is integral iff the subalgebra it generates is finite as an $A$-module. The decisive step is to use the universal property, exactness criterion, localization test, or finite-generation mechanism appropriate to the algebraic structure.
\end{example}

\begin{example}[Finite algebras diagnostic]\label{ex:v19-worked-04}
Give a rigorous worked analysis of Finite algebras. State the ring, module, finiteness, localization, or homological hypotheses and derive the central conclusion. Every finite $A$-algebra is integral over $A$. The decisive step is to use the universal property, exactness criterion, localization test, or finite-generation mechanism appropriate to the algebraic structure.
\end{example}

\section{Solved dossiers}
Each dossier is a canonical solved problem. The provenance ledger records which retained corpus rules guided its topic and which dossiers were added freshly to complete the chapter.

\begin{problem}[Integral elements diagnostic]\label{prob:v19-dossier-01}
Give a rigorous worked analysis of Integral elements. State the ring, module, finiteness, localization, or homological hypotheses and derive the central conclusion.
\end{problem}
\begin{solution}
An element $b$ of an $A$-algebra $B$ is integral over $A$ if it satisfies a monic polynomial with coefficients in $A$. The decisive step is to use the universal property, exactness criterion, localization test, or finite-generation mechanism appropriate to the algebraic structure.
\end{solution}

\begin{problem}[Integral extensions diagnostic]\label{prob:v19-dossier-02}
Give a rigorous worked analysis of Integral extensions. State the ring, module, finiteness, localization, or homological hypotheses and derive the central conclusion.
\end{problem}
\begin{solution}
An extension $A\subset B$ is integral when every element of $B$ is integral over $A$. The decisive step is to use the universal property, exactness criterion, localization test, or finite-generation mechanism appropriate to the algebraic structure.
\end{solution}

\begin{problem}[Finite module criterion diagnostic]\label{prob:v19-dossier-03}
Give a rigorous worked analysis of Finite module criterion. State the ring, module, finiteness, localization, or homological hypotheses and derive the central conclusion.
\end{problem}
\begin{solution}
An element is integral iff the subalgebra it generates is finite as an $A$-module. The decisive step is to use the universal property, exactness criterion, localization test, or finite-generation mechanism appropriate to the algebraic structure.
\end{solution}

\begin{problem}[Finite algebras diagnostic]\label{prob:v19-dossier-04}
Give a rigorous worked analysis of Finite algebras. State the ring, module, finiteness, localization, or homological hypotheses and derive the central conclusion.
\end{problem}
\begin{solution}
Every finite $A$-algebra is integral over $A$. The decisive step is to use the universal property, exactness criterion, localization test, or finite-generation mechanism appropriate to the algebraic structure.
\end{solution}

\begin{problem}[Transitivity diagnostic]\label{prob:v19-dossier-05}
Give a rigorous worked analysis of Transitivity. State the ring, module, finiteness, localization, or homological hypotheses and derive the central conclusion.
\end{problem}
\begin{solution}
Integrality is transitive through towers of ring extensions. The decisive step is to use the universal property, exactness criterion, localization test, or finite-generation mechanism appropriate to the algebraic structure.
\end{solution}

\begin{problem}[Integral elements form a ring diagnostic]\label{prob:v19-dossier-06}
Give a rigorous worked analysis of Integral elements form a ring. State the ring, module, finiteness, localization, or homological hypotheses and derive the central conclusion.
\end{problem}
\begin{solution}
The set of elements of $B$ integral over $A$ is a subring. The decisive step is to use the universal property, exactness criterion, localization test, or finite-generation mechanism appropriate to the algebraic structure.
\end{solution}

\begin{problem}[Determinant trick diagnostic]\label{prob:v19-dossier-07}
Give a rigorous worked analysis of Determinant trick. State the ring, module, finiteness, localization, or homological hypotheses and derive the central conclusion.
\end{problem}
\begin{solution}
Finite-generation plus stability under multiplication produces monic equations by a determinant argument. The decisive step is to use the universal property, exactness criterion, localization test, or finite-generation mechanism appropriate to the algebraic structure.
\end{solution}

\begin{problem}[Prime behavior preview diagnostic]\label{prob:v19-dossier-08}
Give a rigorous worked analysis of Prime behavior preview. State the ring, module, finiteness, localization, or homological hypotheses and derive the central conclusion.
\end{problem}
\begin{solution}
Integral extensions satisfy lying-over and contraction properties for prime ideals. The decisive step is to use the universal property, exactness criterion, localization test, or finite-generation mechanism appropriate to the algebraic structure.
\end{solution}

\begin{problem}[Finite module criterion proof dossier]\label{prob:v19-dossier-09}
Prove the following structural statement and identify the hypothesis that makes the conclusion work: For $b\in B$, $b$ is integral over $A$ iff $A[b]$ is finitely generated as an $A$-module.
\end{problem}
\begin{solution}
A monic equation expresses every high power of $b$ in terms of lower powers. Conversely, if finitely many powers generate an $A$-module stable under multiplication by $b$, the determinant trick gives a monic polynomial annihilating $b$. This proof isolates the reusable algebraic mechanism behind the result.
\end{solution}

\begin{problem}[Finite algebra implies integral proof dossier]\label{prob:v19-dossier-10}
Prove the following structural statement and identify the hypothesis that makes the conclusion work: If $B$ is a finite $A$-module and an $A$-algebra, then every $b\in B$ is integral over $A$.
\end{problem}
\begin{solution}
Multiplication by $b$ is an $A$-linear endomorphism of the finite module $B$. Apply the determinant trick or Cayley--Hamilton to obtain a monic polynomial with coefficients in $A$ annihilating $b$. This proof isolates the reusable algebraic mechanism behind the result.
\end{solution}

\begin{problem}[Transitivity of integrality proof dossier]\label{prob:v19-dossier-11}
Prove the following structural statement and identify the hypothesis that makes the conclusion work: If $B$ is integral over $A$ and $C$ integral over $B$, then $C$ is integral over $A$.
\end{problem}
\begin{solution}
For $c\in C$, its monic equation uses finitely many coefficients from $B$. Those coefficients generate a finite $A$-algebra because they are integral; adjoining $c$ is finite over that finite algebra, hence finite over $A$, so $c$ is integral over $A$. This proof isolates the reusable algebraic mechanism behind the result.
\end{solution}

\begin{problem}[Chapter synthesis]\label{prob:v19-dossier-12}
Connect the definitions and principal structural theorems of this chapter into one reusable commutative-algebra strategy.
\end{problem}
\begin{solution}
Integral dependence generalizes algebraicity from field extensions to ring extensions. An element is integral when it satisfies a monic polynomial, and integrality is stable under finite algebraic constructions. The recurring strategy is to encode the algebraic object by kernels, quotients, localization, finite presentation, or a resolution, apply the corresponding universal or exactness principle, and then recover the desired global or local conclusion.
\end{solution}

\section{Exercises with complete solutions}

\begin{exercise}\label{exr:v19-01}
State and apply the central principle behind Integral elements. Give one concrete algebraic consequence.
\end{exercise}
\begin{hint}
Integral means satisfying a monic polynomial over A; monicity is essential, so do not use an arbitrary algebraic relation.
\end{hint}
\begin{solution}
An element $b$ of an $A$-algebra $B$ is integral over $A$ if it satisfies a monic polynomial with coefficients in $A$. The same principle can then be reused in later localization, support, base-change, resolution, Tor, or Ext arguments.
\end{solution}

\begin{exercise}\label{exr:v19-02}
State and apply the central principle behind Integral extensions. Give one concrete algebraic consequence.
\end{exercise}
\begin{hint}
To prove closure under addition/multiplication, place the elements inside a common finite A-module algebra and use the finite-module criterion.
\end{hint}
\begin{solution}
An extension $A\subset B$ is integral when every element of $B$ is integral over $A$. The same principle can then be reused in later localization, support, base-change, resolution, Tor, or Ext arguments.
\end{solution}

\begin{exercise}\label{exr:v19-03}
State and apply the central principle behind Finite module criterion. Give one concrete algebraic consequence.
\end{exercise}
\begin{hint}
The determinant trick starts from a finite module stable under multiplication by x and produces a monic characteristic relation.
\end{hint}
\begin{solution}
An element is integral iff the subalgebra it generates is finite as an $A$-module. The same principle can then be reused in later localization, support, base-change, resolution, Tor, or Ext arguments.
\end{solution}

\begin{exercise}\label{exr:v19-04}
State and apply the central principle behind Finite algebras. Give one concrete algebraic consequence.
\end{exercise}
\begin{hint}
Integrality descends to quotients by reducing the monic polynomial modulo the ideal.
\end{hint}
\begin{solution}
Every finite $A$-algebra is integral over $A$. The same principle can then be reused in later localization, support, base-change, resolution, Tor, or Ext arguments.
\end{solution}

\begin{exercise}\label{exr:v19-05}
State and apply the central principle behind Transitivity. Give one concrete algebraic consequence.
\end{exercise}
\begin{hint}
Integrality localizes because the same monic relation survives after localizing coefficients.
\end{hint}
\begin{solution}
Integrality is transitive through towers of ring extensions. The same principle can then be reused in later localization, support, base-change, resolution, Tor, or Ext arguments.
\end{solution}

\begin{exercise}\label{exr:v19-06}
State and apply the central principle behind Integral elements form a ring. Give one concrete algebraic consequence.
\end{exercise}
\begin{hint}
For prime behavior, extend/contraction arguments use integrality together with the lying-over mechanism; identify the exact theorem before applying it.
\end{hint}
\begin{solution}
The set of elements of $B$ integral over $A$ is a subring. The same principle can then be reused in later localization, support, base-change, resolution, Tor, or Ext arguments.
\end{solution}

\begin{exercise}\label{exr:v19-07}
State and apply the central principle behind Determinant trick. Give one concrete algebraic consequence.
\end{exercise}
\begin{hint}
Every element of a finite A-algebra is integral over A; use Cayley--Hamilton on multiplication by that element.
\end{hint}
\begin{solution}
Finite-generation plus stability under multiplication produces monic equations by a determinant argument. The same principle can then be reused in later localization, support, base-change, resolution, Tor, or Ext arguments.
\end{solution}

\begin{exercise}\label{exr:v19-08}
State and apply the central principle behind Prime behavior preview. Give one concrete algebraic consequence.
\end{exercise}
\begin{hint}
Separate 'algebraic over the fraction field' from 'integral over A': denominators in a polynomial relation matter.
\end{hint}
\begin{solution}
Integral extensions satisfy lying-over and contraction properties for prime ideals. The same principle can then be reused in later localization, support, base-change, resolution, Tor, or Ext arguments.
\end{solution}

\section*{Chapter summary}
The chapter now contributes a canonical solved layer to the progression from rings and modules through localization, Noetherian methods, integral dependence, and derived functors.
% BEGIN VOL05-EXPANSION V19-exercises-01
\section{Graded supplementary exercises}
These exercises extend the protected chapter with a balanced set of computations, proofs, hypothesis tests, applications, and challenges.

\subsection*{Standard computations}

\begin{exercise}[Square root]\label{exr:v19-expand-01}
Is $\sqrt2$ integral over $\mathbb Z$?
\end{exercise}
\begin{hint}
Use monic equations and finite module criteria.
\end{hint}
\begin{solution}
Yes, it satisfies $X^2-2=0$.
\end{solution}

\begin{exercise}[Rational]\label{exr:v19-expand-02}
Is $1/2$ integral over $\mathbb Z$?
\end{exercise}
\begin{hint}
Use monic equations and finite module criteria.
\end{hint}
\begin{solution}
No.
\end{solution}

\begin{exercise}[Cusp parameter]\label{exr:v19-expand-03}
Is $t$ integral over $k[t^2,t^3]$?
\end{exercise}
\begin{hint}
Use monic equations and finite module criteria.
\end{hint}
\begin{solution}
Yes.
\end{solution}

\begin{exercise}[Finite algebra]\label{exr:v19-expand-04}
If $b^3+a_2b^2+a_1b+a_0=0$, which powers generate $A[b]$?
\end{exercise}
\begin{hint}
Use monic equations and finite module criteria.
\end{hint}
\begin{solution}
$1,b,b^2$.
\end{solution}

\begin{exercise}[Transitivity]\label{exr:v19-expand-05}
If $C$ is integral over $B$ and $B$ over $A$, what follows?
\end{exercise}
\begin{hint}
Use monic equations and finite module criteria.
\end{hint}
\begin{solution}
$C$ is integral over $A$.
\end{solution}

\subsection*{Proofs}

\begin{exercise}[Finite module criterion proof]\label{exr:v19-expand-06}
Prove: For $b\in B$, $b$ is integral over $A$ iff $A[b]$ is finitely generated as an $A$-module.
\end{exercise}
\begin{hint}
Start from the theorem hypotheses and identify the decisive algebraic construction.
\end{hint}
\begin{solution}
A monic equation expresses every high power of $b$ in terms of lower powers. Conversely, if finitely many powers generate an $A$-module stable under multiplication by $b$, the determinant trick gives a monic polynomial annihilating $b$.
\end{solution}

\begin{exercise}[Finite algebra implies integral proof]\label{exr:v19-expand-07}
Prove: If $B$ is a finite $A$-module and an $A$-algebra, then every $b\in B$ is integral over $A$.
\end{exercise}
\begin{hint}
Start from the theorem hypotheses and identify the decisive algebraic construction.
\end{hint}
\begin{solution}
Multiplication by $b$ is an $A$-linear endomorphism of the finite module $B$. Apply the determinant trick or Cayley--Hamilton to obtain a monic polynomial with coefficients in $A$ annihilating $b$.
\end{solution}

\begin{exercise}[Transitivity of integrality proof]\label{exr:v19-expand-08}
Prove: If $B$ is integral over $A$ and $C$ integral over $B$, then $C$ is integral over $A$.
\end{exercise}
\begin{hint}
Start from the theorem hypotheses and identify the decisive algebraic construction.
\end{hint}
\begin{solution}
For $c\in C$, its monic equation uses finitely many coefficients from $B$. Those coefficients generate a finite $A$-algebra because they are integral; adjoining $c$ is finite over that finite algebra, hence finite over $A$, so $c$ is integral over $A$.
\end{solution}

\begin{exercise}[Structural synthesis]\label{exr:v19-expand-09}
Prove a corollary that genuinely uses both Finite module criterion and Finite algebra implies integral.
\end{exercise}
\begin{hint}
Apply the first theorem to obtain its structural description, then verify the second theorem's hypotheses.
\end{hint}
\begin{solution}
A correct proof explicitly invokes both results, verifies the common hypotheses, and derives a new consequence rather than merely restating either theorem.
\end{solution}

\subsection*{Counterexamples and hypothesis tests}

\begin{exercise}[Arbitrary fraction]\label{exr:v19-expand-10}
Test the claim: Every element of the fraction field of an integral domain is integral over the domain.
\end{exercise}
\begin{hint}
Use the smallest concrete ring, module, or resolution that can decide the claim.
\end{hint}
\begin{solution}
False: $1/2$ over $\mathbb Z$.
\end{solution}

\begin{exercise}[Monic condition]\label{exr:v19-expand-11}
Test the claim: A root of any polynomial with coefficients in $A$ is automatically integral.
\end{exercise}
\begin{hint}
Use the smallest concrete ring, module, or resolution that can decide the claim.
\end{hint}
\begin{solution}
False; the polynomial must be monic or an equivalent criterion must apply.
\end{solution}

\begin{exercise}[Sum product]\label{exr:v19-expand-12}
Test the claim: The sum and product of integral elements need not be integral.
\end{exercise}
\begin{hint}
Use the smallest concrete ring, module, or resolution that can decide the claim.
\end{hint}
\begin{solution}
False: integral elements form a subring.
\end{solution}

\subsection*{Applications and investigations}

\begin{exercise}[Finite maps]\label{exr:v19-expand-13}
Why does integrality correspond to finite algebra behavior?
\end{exercise}
\begin{hint}
Translate the question into monic equations and finite module criteria.
\end{hint}
\begin{solution}
An integral finite-type algebra is finite as a module.
\end{solution}

\begin{exercise}[Eliminating parameters]\label{exr:v19-expand-14}
Why is $k[t]$ integral over the cusp ring $k[t^2,t^3]$?
\end{exercise}
\begin{hint}
Translate the question into monic equations and finite module criteria.
\end{hint}
\begin{solution}
The missing parameter satisfies a monic relation over the cusp ring.
\end{solution}

\subsection*{Challenge problems}

\begin{exercise}[Local-global synthesis]\label{exr:v19-expand-15}
Connect 'Integral elements' with 'Examples' in one rigorous argument.
\end{exercise}
\begin{hint}
State the relevant ring map, module map, or complex before applying the structural theorem.
\end{hint}
\begin{solution}
The opening idea is: An element $b$ of an $A$-algebra $B$ is integral over $A$ if it satisfies a monic polynomial with coefficients in $A$. The later viewpoint is: Algebraic integers, roots of monic equations, and normalization rings illustrate integral dependence. A complete solution identifies the structural theorem that links these descriptions and checks its hypotheses.
\end{solution}

\begin{exercise}[Hypothesis audit]\label{exr:v19-expand-16}
Choose one theorem from the chapter, remove one essential hypothesis, and exhibit or explain a concrete failure.
\end{exercise}
\begin{hint}
Use one of the chapter's explicit computations or counterexamples as the test case.
\end{hint}
\begin{solution}
The solution must name the removed hypothesis, show exactly which proof step breaks, and give a concrete algebraic object where the claimed conclusion fails.
\end{solution}

% END VOL05-EXPANSION V19-exercises-01
